\documentclass[11pt]{article}
\usepackage[margin=1in]{geometry}
\usepackage{amsmath,amssymb,amsthm,mathtools}
\usepackage[colorlinks=true,linkcolor=blue,citecolor=blue,urlcolor=blue]{hyperref}

\newtheorem{theorem}{Theorem}
\renewcommand{\thetheorem}{\Alph{theorem}}
\newtheorem{lemma}{Lemma}
\newtheorem{proposition}[lemma]{Proposition}
\newtheorem{corollary}[lemma]{Corollary}
\theoremstyle{definition}
\newtheorem{example}[lemma]{Example}
\theoremstyle{remark}
\newtheorem{remark}[lemma]{Remark}

\newcommand{\R}{\mathbb{R}}
\newcommand{\Sph}{\mathbb{S}}
\newcommand{\Tr}{\operatorname{Tr}}
\newcommand{\dd}{\,\mathrm{d}}
\newcommand{\FC}[1]{\mathcal{F}_{#1}}
\newcommand{\Ct}{\widetilde{C}_t}
\newcommand{\taut}{\widetilde{\tau}_t}
\newcommand{\ip}[2]{\langle #1,#2\rangle}
\newcommand{\hatf}[1]{\widehat{#1}}
\newcommand{\tp}{\top}

\title{The Lyapunov-like condition for Koopman invariance of Gaussian RKHSs\\
is not necessary: an exact characterization}
\author{}
\date{September 2026}

\begin{document}
\maketitle

\begin{abstract}
We settle the question left open by Philipp, Schaller, Worthmann, Peitz and N\"uske \cite{PSWPN}.
For the Ornstein--Uhlenbeck equation $\mathrm{d}X_t=AX_t\,\mathrm{d}t+B\,\mathrm{d}W_t$ and the Gaussian
kernel $k^C(x,y)=\exp(-\|C^{-1}(x-y)\|^2)$, the reproducing kernel Hilbert space $H_C$ is invariant under
every Koopman operator $K^t$, $t\ge 0$, \emph{if and only if}
\[
\tfrac12\bigl(AC^2+C^2A^\tp\bigr)\le 2BB^\tp .
\]
The Lyapunov-like condition $\tfrac12(AC^2+C^2A^\tp)\le BB^\tp$ of \cite[Theorem~2.1]{PSWPN} is therefore
sufficient but \emph{not} necessary; we give explicit counterexamples that satisfy all standing assumptions
of \cite{PSWPN}. At a single time $t$, $K^tH_C\subset H_C$ holds if and only if
$e^{At}C^2e^{A^\tp t}\le C^2+4\Sigma(t)$. Moreover, $K^t$ maps $H_C$ bijectively onto the Gaussian RKHS
$H_{\widetilde C_t}$ with $\widetilde C_t^{\,2}=e^{-At}(C^2+4\Sigma(t))e^{-A^\tp t}$. For $t>0$ this space is strictly
smaller than the space $H_{C_t}$, $C_t^2=e^{-At}(C^2+2\Sigma(t))e^{-A^\tp t}$, of \cite[Proposition~4.1]{PSWPN}; the discrepancy
$2\Sigma(t)$ versus $4\Sigma(t)$ is exactly why the sufficient condition of \cite{PSWPN} is not sharp. Whenever invariance holds,
$\|K^t\|_{H_C\to H_C}=e^{t\Tr(-A)/2}$, so the bound of \cite[Theorem~2.1]{PSWPN} is attained. For Hurwitz
$A$ the criterion says that the ellipsoid of $C^2+4\Sigma$ is forward invariant under $\dot x=Ax$.
\end{abstract}

\tableofcontents

%======================================================================
\section{The problem and the answer}\label{sec:problem}
%======================================================================

\paragraph{Setting of \cite{PSWPN}.}
Let $A\in\R^{d\times d}$, $B\in\R^{d\times m}$, let $W_t$ be an $m$-dimensional Brownian motion, and consider
\begin{equation}\label{eq:sde}
\mathrm{d}X_t=AX_t\,\mathrm{d}t+B\,\mathrm{d}W_t .
\end{equation}
Put
\begin{equation}\label{eq:Sigmat}
\Sigma(t):=\int_0^t e^{As}BB^\tp e^{A^\tp s}\dd s\qquad(t\ge 0).
\end{equation}
In \cite{PSWPN} it is assumed that $m\le d$, that $(A,B)$ is controllable and that $A$ is Hurwitz; then
$\Sigma:=\lim_{t\to\infty}\Sigma(t)$ exists, is positive definite, solves $A\Sigma+\Sigma A^\tp=-BB^\tp$, and
$\mu:=\mathcal N(0,\Sigma)$ is the stationary distribution. The Koopman operator is
$(K^tf)(x):=\mathbb E[f(X_t)\,|\,X_0=x]$. For a symmetric positive definite $C\in\R^{d\times d}$,
\[
k^C(x,y):=\exp\bigl(-(x-y)^\tp C^{-2}(x-y)\bigr)=\exp\bigl(-\|C^{-1}(x-y)\|^2\bigr),\qquad x,y\in\R^d,
\]
$H_C$ denotes the RKHS of $k^C$ on $\R^d$, $\|\cdot\|_C$ its norm, and $k^C_x:=k^C(x,\cdot)$. Inequalities between
symmetric matrices refer to the Loewner order.

\paragraph{The result of \cite{PSWPN} and the open problem.}
\cite[Theorem~2.1]{PSWPN}: if
\begin{equation}\tag{2.3}\label{eq:23}
\tfrac12\bigl(AC^2+C^2A^\tp\bigr)\le BB^\tp ,
\end{equation}
then $K^tH_C\subset H_C$ for all $t\ge0$ and $\|K^tf\|_C\le e^{t\Tr(-A)/2}\|f\|_C$. The authors ask
(\cite[Abstract, Remark~2.2(c), Remark~4.2]{PSWPN}) whether \eqref{eq:23} is also \emph{necessary} for
$K^tH_C\subset H_C$, $t\ge 0$. In Remark~4.2 they point out that the converse would follow if one could show
that $K^tH_C\subset H_C$ implies $H_{C_t}\hookrightarrow H_C$, where
$C_t=\bigl(e^{-At}(C^2+2\Sigma(t))e^{-A^\tp t}\bigr)^{1/2}$ is the matrix of \cite[(4.1)]{PSWPN}.

\paragraph{Answer.}
\emph{No}: \eqref{eq:23} is not necessary. The necessary and sufficient condition is
\begin{equation}\tag{2.3$^\sharp$}\label{eq:23sharp}
\tfrac12\bigl(AC^2+C^2A^\tp\bigr)\le 2BB^\tp ,
\end{equation}
that is, \eqref{eq:23} with $BB^\tp$ replaced by $2BB^\tp$; equivalently, \eqref{eq:23} holds for $C/\sqrt2$ in place
of $C$ (Theorem~\ref{thm:C}, Corollary~\ref{cor:answer}). Examples~\ref{ex:beta} and~\ref{ex:scalar} satisfy
\eqref{eq:23sharp} but violate \eqref{eq:23}; they fulfil all standing assumptions of \cite{PSWPN}. The remaining
results describe the situation completely:
\begin{itemize}
\item \emph{fixed time} (Theorem~\ref{thm:A}): $K^tH_C\subset H_C$ iff $e^{At}C^2e^{A^\tp t}\le C^2+4\Sigma(t)$, and then
$\|K^t\|_{H_C\to H_C}=e^{t\Tr(-A)/2}$;
\item \emph{range} (Theorem~\ref{thm:B}): $K^t$ maps $H_C$ bijectively onto $H_{\Ct}$,
$\Ct^{\,2}=e^{-At}(C^2+4\Sigma(t))e^{-A^\tp t}$, and $H_{\Ct}\subsetneq H_{C_t}$ whenever $\Sigma(t)\neq0$. The
implication proposed in \cite[Remark~4.2]{PSWPN} is false in general (Corollary~\ref{cor:remark42});
\item \emph{geometry} (Theorem~\ref{thm:D}, $A$ Hurwitz): $K^tH_C\subset H_C$ iff $e^{At}$ is non-expansive in the
norm $x\mapsto(x^\tp(C^2+4\Sigma)^{-1}x)^{1/2}$. Condition~\eqref{eq:23sharp} holds iff
$A(C^2+4\Sigma)+(C^2+4\Sigma)A^\tp\le0$, whereas \eqref{eq:23} holds iff $A(C^2+2\Sigma)+(C^2+2\Sigma)A^\tp\le 0$.
\end{itemize}

\paragraph{Scope and conventions.}
For bounded Borel $f$ we use the pointwise definition $(K^tf)(x):=\mathbb E[f(X_t)\,|\,X_0=x]$. Elements of $H_C$ are
bounded and continuous, and $K^tf$ is then continuous (Lemma~\ref{lem:Kt}). Under the standing assumptions of
\cite{PSWPN}, $K^t$ is considered on $L^2(\mu)$ and $H_C$ is identified with a subspace of $L^2(\mu)$
\cite[Lemma~3.1]{PSWPN}. Since $\mu$ has a positive density and elements of $H_C$ are continuous, the
$L^2(\mu)$-class of $K^tf$ contains an element of $H_C$ if and only if the continuous function $K^tf$ itself
belongs to $H_C$. Hence the two readings of ``$K^tH_C\subset H_C$'' coincide. \emph{Sections~\ref{sec:fourier}--\ref{sec:main}
use neither controllability, nor the Hurwitz property, nor $m\le d$}: they hold for arbitrary $A\in\R^{d\times d}$, $B\in\R^{d\times m}$.
The Hurwitz property enters only Theorem~\ref{thm:D}. Everything in Sections~\ref{sec:fourier}--\ref{sec:remarks}
is proved analytically, except the paragraph marked \emph{Numerical illustration} in Example~\ref{ex:beta}.
Section~\ref{sec:numerics} reports floating-point cross-checks only, and no proof relies on them.

%======================================================================
\section{Fourier description of \texorpdfstring{$H_C$}{H\_C}}\label{sec:fourier}
%======================================================================

\paragraph{Conventions.}
For $g\in L^1(\R^d)$ let $\hatf g(\xi):=\int_{\R^d}g(x)e^{-i\xi\cdot x}\dd x$. The map $g\mapsto\hatf g$ extends to $L^2(\R^d)$
with $\|\hatf g\|_{L^2}^2=(2\pi)^d\|g\|_{L^2}^2$ (Plancherel). For $F\in L^1(\R^d)\cap L^2(\R^d)$ put
$F^\vee(x):=(2\pi)^{-d}\int F(\xi)e^{i\xi\cdot x}\dd\xi$. Then $F^\vee$ is continuous, $F^\vee\in L^2(\R^d)$ and
$\hatf{F^\vee}=F$. If $g\in L^2(\R^d)\cap C(\R^d)$ and $\hatf g\in L^1(\R^d)$, then $g=(\hatf g)^\vee$ everywhere:
the two functions are continuous and agree almost everywhere. A measurable $F$ is \emph{Hermitian} if
$F(-\xi)=\overline{F(\xi)}$ for a.e.\ $\xi$. For real-valued $g$, $\hatf g$ is Hermitian, and for Hermitian
$F\in L^1\cap L^2$, $F^\vee$ is real-valued. For symmetric positive definite $P\in\R^{d\times d}$ and
$\xi\in\R^d$,
\begin{equation}\label{eq:gauss}
\int_{\R^d}\exp\bigl(-x^\tp Px-i\xi\cdot x\bigr)\dd x=\pi^{d/2}(\det P)^{-1/2}\exp\bigl(-\tfrac14\xi^\tp P^{-1}\xi\bigr),
\end{equation}
which follows by orthogonal diagonalisation of $P$ from $\int_\R e^{-px^2-i\xi x}\dd x=\sqrt{\pi/p}\,e^{-\xi^2/(4p)}$.

For symmetric positive definite $C$ define
\begin{equation}\label{eq:wC}
w_C(\xi):=\exp\bigl(\tfrac14\xi^\tp C^2\xi\bigr),\qquad \kappa_C:=\pi^{d/2}\det C,\qquad
\FC{C}:=\Bigl\{F:\R^d\to\mathbb C\ \text{measurable, Hermitian},\ \int|F|^2w_C\dd\xi<\infty\Bigr\}.
\end{equation}
Applying \eqref{eq:gauss} with $P=C^{-2}$ and a translation,
\begin{equation}\label{eq:kernelFT}
\hatf{k^C_y}(\xi)=\kappa_C\,e^{-i\xi\cdot y}\,w_C(\xi)^{-1},\qquad y,\xi\in\R^d ,
\end{equation}
and, by \eqref{eq:gauss} with $\xi=0$ and $P=\frac14C^2$,
\begin{equation}\label{eq:intw}
\int_{\R^d}w_C(\xi)^{-1}\dd\xi=\frac{(4\pi)^{d/2}}{\det C}<\infty .
\end{equation}

\begin{lemma}[Fourier description of $H_C$]\label{lem:fourier}
Let $C\in\R^{d\times d}$ be symmetric positive definite.
\begin{enumerate}
\item[(a)] Every $F\in\FC{C}$ belongs to $L^1(\R^d)\cap L^2(\R^d)$, and
$\|F\|_{L^1}\le\bigl((4\pi)^{d/2}/\det C\bigr)^{1/2}\bigl(\int|F|^2w_C\bigr)^{1/2}$.
\item[(b)] $H_C=\{f\in C(\R^d)\cap L^2(\R^d)\ \text{real-valued}:\ \hatf f\in\FC{C}\}$, and for $f,g\in H_C$
\begin{equation}\label{eq:normC}
\ip fg_C=(2\pi)^{-d}\kappa_C^{-1}\int_{\R^d}\hatf f(\xi)\,\overline{\hatf g(\xi)}\,w_C(\xi)\dd\xi .
\end{equation}
\item[(c)] $f\mapsto\hatf f$ is a bijection of $H_C$ onto $\FC{C}$ with inverse $F\mapsto F^\vee$.
\end{enumerate}
\end{lemma}

\begin{proof}
(a) By the Cauchy--Schwarz inequality and \eqref{eq:intw},
$\int|F|=\int|F|w_C^{1/2}\,w_C^{-1/2}\le(\int|F|^2w_C)^{1/2}(\int w_C^{-1})^{1/2}$. Since $w_C\ge1$, also
$\int|F|^2\le\int|F|^2w_C$.

(b) Let $\mathcal H$ be the set on the right-hand side, equipped with the form $\ip fg_{\mathcal H}$ given by the
right-hand side of \eqref{eq:normC}.

\emph{$\mathcal H$ is a real inner product space.} For real $f,g$, the functions $\hatf f$ and $\hatf g$ are Hermitian. Because $w_C$
is even, the integrand at $-\xi$ is the complex conjugate of the integrand at $\xi$, so the integral is real.
The form is bilinear, symmetric and nonnegative. If $\ip ff_{\mathcal H}=0$ then $\hatf f=0$ a.e., hence $f=0$
in $L^2$, hence $f=0$ by continuity.

\emph{Inversion.} For $f\in\mathcal H$ we have $\hatf f\in L^1$ by (a), hence
\begin{equation}\label{eq:inversion}
f(x)=(2\pi)^{-d}\int\hatf f(\xi)e^{i\xi\cdot x}\dd\xi\qquad\text{for every }x\in\R^d .
\end{equation}

\emph{Completeness.} Let $(f_n)$ be Cauchy in $\mathcal H$. Then $(\hatf{f_n})$ is Cauchy in $L^2(w_C\dd\xi)$ and
converges there to some $F$. A subsequence converges a.e., so $F$ is Hermitian, and $F\in\FC{C}$. By (a),
$F\in L^1\cap L^2$. Then $f:=F^\vee$ is continuous, real-valued, lies in $L^2$ and satisfies $\hatf f=F$. Hence $f\in\mathcal H$ and
$\|f_n-f\|_{\mathcal H}^2=(2\pi)^{-d}\kappa_C^{-1}\int|\hatf{f_n}-F|^2w_C\to0$.

\emph{Kernel sections.} By \eqref{eq:kernelFT} and \eqref{eq:intw}, $\hatf{k^C_y}$ is Hermitian with
$\int|\hatf{k^C_y}|^2w_C=\kappa_C^2\int w_C^{-1}<\infty$. Since $k^C_y$ is real, continuous and in $L^2$, we have
$k^C_y\in\mathcal H$.

\emph{Reproducing property.} For $f\in\mathcal H$ and $y\in\R^d$, \eqref{eq:kernelFT} and \eqref{eq:inversion} give
\[
\ip f{k^C_y}_{\mathcal H}=(2\pi)^{-d}\kappa_C^{-1}\int\hatf f(\xi)\,\kappa_C e^{i\xi\cdot y}w_C(\xi)^{-1}w_C(\xi)\dd\xi
=(2\pi)^{-d}\int\hatf f(\xi)e^{i\xi\cdot y}\dd\xi=f(y).
\]

\emph{Identification with $H_C$.} Let $H_0:=\operatorname{span}\{k^C_y:y\in\R^d\}$. By the reproducing property, $H_0$
is dense in $\mathcal H$ (if $f\perp H_0$ then $f(y)=\ip f{k^C_y}_{\mathcal H}=0$ for all $y$). It is also dense
in $H_C$ by the same argument. On $H_0$ both norms coincide:
$\|\sum_j\alpha_jk^C_{y_j}\|^2=\sum_{i,j}\alpha_i\alpha_jk^C(y_i,y_j)$. In both spaces norm convergence implies
pointwise convergence, because $|h(x)|\le\|h\|\,k^C(x,x)^{1/2}=\|h\|$. Given $f\in\mathcal H$, choose $f_n\in H_0$
with $f_n\to f$ in $\mathcal H$. Then $(f_n)$ is Cauchy in $H_C$, so $f_n\to g$ in $H_C$ for some $g$, and the pointwise
limits give $g=f$. Thus $f\in H_C$ with $\|f\|_C=\lim\|f_n\|_C=\lim\|f_n\|_{\mathcal H}=\|f\|_{\mathcal H}$.
Exchanging the roles of $\mathcal H$ and $H_C$ gives $H_C\subset\mathcal H$. Hence $\mathcal H=H_C$ isometrically,
and \eqref{eq:normC} follows by polarisation. (This is the uniqueness part of the Moore--Aronszajn theorem,
cf.\ \cite{Aronszajn,PaulsenRaghupathi}.)

(c) By (b), $\hatf f\in\FC{C}$ for $f\in H_C$. The map is injective, since $\hatf f$ determines the continuous
$L^2$-function $f$. It is surjective: for $F\in\FC{C}$, (a) gives $F\in L^1\cap L^2$, so $F^\vee$ is real, continuous,
in $L^2$, with $\hatf{F^\vee}=F\in\FC{C}$, and $F^\vee\in H_C$ by (b).
\end{proof}

Lemma~\ref{lem:fourier} is the Gaussian case of the Fourier description of native spaces of
translation-invariant kernels, cf.\ \cite[Theorem~10.12]{Wendland}. The proof above is self-contained. A different,
power-series description of Gaussian RKHSs is given in \cite{SHS}.

\begin{lemma}[Exponential multipliers]\label{lem:multiplier}
Let $C$ be symmetric positive definite and let $S\in\R^{d\times d}$ be symmetric. Then
\[
\int_{\R^d}|F(\xi)|^2\,w_C(\xi)\,\exp\bigl(\tfrac14\xi^\tp S\xi\bigr)\dd\xi<\infty\qquad\text{for every }F\in\FC{C}
\]
if and only if $S\le0$. If $S\le 0$ fails, there is a real-valued even function $F_S\in\FC{C}$, supported in an open
double cone, for which the integral is $+\infty$.
\end{lemma}

\begin{proof}
If $S\le0$, then $\exp(\frac14\xi^\tp S\xi)\le1$. Conversely, suppose that $S\le0$ fails, and let $\lambda>0$ be an
eigenvalue of $S$ with unit eigenvector $v$. The set $U:=\{\theta\in\Sph^{d-1}:\theta^\tp S\theta>\lambda/2\}$ is relatively open in $\Sph^{d-1}$, satisfies $U=-U$, and contains $v$.
The set $\Gamma:=\{r\theta:\ r>0,\ \theta\in U\}$ is an open double cone, and $\xi^\tp S\xi>\frac\lambda2|\xi|^2$ for
$\xi\in\Gamma$. Define
\[
F_S(\xi):=w_C(\xi)^{-1/2}\,(1+|\xi|^2)^{-d/2}\,\mathbf 1_\Gamma(\xi).
\]
$F_S$ is real and even (because $\Gamma=-\Gamma$ and $w_C$ is even), hence Hermitian. Moreover
$\int|F_S|^2w_C=\int_\Gamma(1+|\xi|^2)^{-d}\dd\xi\le\int_{\R^d}(1+|\xi|^2)^{-d}\dd\xi<\infty$, since the radial integrand
$r^{d-1}(1+r^2)^{-d}$ is $O(r^{-d-1})$ at infinity. So $F_S\in\FC{C}$. In polar coordinates, with $|U|>0$ the surface
measure of $U$ (counting measure if $d=1$),
\[
\int|F_S|^2w_C\,e^{\frac14\xi^\tp S\xi}\dd\xi\ \ge\ \int_\Gamma(1+|\xi|^2)^{-d}e^{\lambda|\xi|^2/8}\dd\xi
=|U|\int_0^\infty r^{d-1}(1+r^2)^{-d}e^{\lambda r^2/8}\dd r=+\infty. \qedhere
\]
\end{proof}

\begin{corollary}[Fourier proof of {\cite[Proposition~3.3]{PSWPN}}]\label{cor:inclusion}
For symmetric positive definite $C_1,C_2$ we have $H_{C_1}\subset H_{C_2}$ if and only if $C_1^2\ge C_2^2$. In
this case $\|f\|_{C_2}\le(\det C_1/\det C_2)^{1/2}\|f\|_{C_1}$ for $f\in H_{C_1}$.
\end{corollary}

\begin{proof}
By Lemma~\ref{lem:fourier}(b),(c), $H_{C_1}\subset H_{C_2}$ holds iff $\int|F|^2w_{C_2}<\infty$ for every
$F\in\FC{C_1}$. Since $w_{C_2}=w_{C_1}\exp(\frac14\xi^\tp(C_2^2-C_1^2)\xi)$, Lemma~\ref{lem:multiplier} with
$S=C_2^2-C_1^2$ gives the equivalence. If $C_1^2\ge C_2^2$, then $w_{C_2}\le w_{C_1}$, and by \eqref{eq:normC}
$\|f\|_{C_2}^2\le(\kappa_{C_1}/\kappa_{C_2})\|f\|_{C_1}^2=(\det C_1/\det C_2)\|f\|_{C_1}^2$.
\end{proof}

%======================================================================
\section{The Koopman operators in Fourier variables}\label{sec:koopman}
%======================================================================

Let $\gamma_t:=\mathcal N(0,\Sigma(t))$, a centred Gaussian measure on $\R^d$ that may be degenerate
($\gamma_0=\delta_0$). The solution of \eqref{eq:sde} with $X_0=x$ is $X_t=e^{At}x+Z_t$, where
$Z_t=\int_0^te^{A(t-s)}B\dd W_s$ is centred Gaussian with covariance $\int_0^te^{A(t-s)}BB^\tp e^{A^\tp(t-s)}\dd s=\Sigma(t)$
(It\^o isometry). Hence, for bounded Borel $f$,
\begin{equation}\label{eq:Ktconv}
(K^tf)(x)=\int_{\R^d}f(e^{At}x+y)\,\gamma_t(\mathrm dy),\qquad x\in\R^d .
\end{equation}
Under the standing assumptions this is \cite[(2.2)]{PSWPN}. Every $f\in H_C$ is bounded and continuous:
$|f(x)|=|\ip f{k^C_x}_C|\le\|f\|_C\,k^C(x,x)^{1/2}=\|f\|_C$.

\begin{lemma}\label{lem:Kt}
Let $f\in H_C$ and $t\ge0$. Then $K^tf$ is real-valued, continuous, belongs to $L^2(\R^d)$, and for a.e.\ $\xi$
\begin{equation}\label{eq:KtFT}
\hatf{K^tf}(\xi)=e^{-t\Tr A}\,\hatf f\bigl(e^{-A^\tp t}\xi\bigr)\,
\exp\Bigl(-\tfrac12\,\xi^\tp e^{-At}\Sigma(t)e^{-A^\tp t}\xi\Bigr).
\end{equation}
\end{lemma}

\begin{proof}
Continuity follows from \eqref{eq:Ktconv} by dominated convergence, since $|f|\le\|f\|_C$ and $f$ is continuous.
For $g\in L^2(\R^d)$ put $(J_tg)(z):=\int g(z+y)\gamma_t(\mathrm dy)$. By Minkowski's integral inequality applied to
$|g|$, the integral converges absolutely for a.e.\ $z$, and $\|J_tg\|_{L^2}\le\int\|g(\cdot+y)\|_{L^2}\gamma_t(\mathrm dy)=\|g\|_{L^2}$.
For $g\in L^1\cap L^2$, Fubini's theorem and the characteristic function of $\gamma_t$ give
\[
\hatf{J_tg}(\xi)=\int\!\!\int g(z+y)e^{-i\xi\cdot z}\dd z\,\gamma_t(\mathrm dy)=\hatf g(\xi)\int e^{i\xi\cdot y}\gamma_t(\mathrm dy)
=\hatf g(\xi)\,e^{-\frac12\xi^\tp\Sigma(t)\xi}.
\]
Both sides are bounded linear maps from $L^2$ to $L^2$ in $g$, so the identity extends from the dense subspace
$L^1\cap L^2$ to all of $L^2$. For $f\in H_C$ the integral defining $J_tf$ converges everywhere, and by
\eqref{eq:Ktconv} $K^tf=(J_tf)\circ e^{At}$. Finally, for $u\in L^2$ and invertible $E$, $u\circ E\in L^2$ and
$\hatf{u\circ E}(\xi)=|\det E|^{-1}\hatf u(E^{-\tp}\xi)$: substitute $x=E^{-1}z$ for $u\in L^1\cap L^2$ and extend by density.
With $E=e^{At}$, $\det E=e^{t\Tr A}$ and $E^{-\tp}=e^{-A^\tp t}$ this yields \eqref{eq:KtFT}, because
$(e^{-A^\tp t}\xi)^\tp\Sigma(t)(e^{-A^\tp t}\xi)=\xi^\tp e^{-At}\Sigma(t)e^{-A^\tp t}\xi$.
\end{proof}

Define the symmetric matrices
\begin{equation}\label{eq:Dt}
D_t:=C^2+4\Sigma(t)-e^{At}C^2e^{A^\tp t}\quad(t\ge0),\qquad
\Lambda_C:=4BB^\tp-\bigl(AC^2+C^2A^\tp\bigr).
\end{equation}
Note that \eqref{eq:23sharp} is equivalent to $\Lambda_C\ge0$, and \eqref{eq:23} to $2BB^\tp-(AC^2+C^2A^\tp)\ge0$.

\begin{proposition}[Norm identity]\label{prop:norm}
For $f\in H_C$ and $t\ge0$,
\begin{equation}\label{eq:normid}
\int_{\R^d}\bigl|\hatf{K^tf}(\xi)\bigr|^2w_C(\xi)\dd\xi
=e^{-t\Tr A}\int_{\R^d}|\hatf f(\eta)|^2\,w_C(\eta)\,\exp\bigl(-\tfrac14\eta^\tp D_t\eta\bigr)\dd\eta\ \in[0,\infty].
\end{equation}
\end{proposition}

\begin{proof}
By \eqref{eq:KtFT}, the left-hand side equals
$e^{-2t\Tr A}\int|\hatf f(e^{-A^\tp t}\xi)|^2\exp\bigl(-(e^{-A^\tp t}\xi)^\tp\Sigma(t)(e^{-A^\tp t}\xi)+\frac14\xi^\tp C^2\xi\bigr)\dd\xi$.
Substituting $\xi=e^{A^\tp t}\eta$, $\mathrm d\xi=e^{t\Tr A}\mathrm d\eta$, it becomes
\[
e^{-t\Tr A}\int|\hatf f(\eta)|^2\exp\Bigl(-\eta^\tp\Sigma(t)\eta+\tfrac14\eta^\tp e^{At}C^2e^{A^\tp t}\eta\Bigr)\dd\eta .
\]
The claim follows from
$-\eta^\tp\Sigma(t)\eta+\frac14\eta^\tp e^{At}C^2e^{A^\tp t}\eta=\frac14\eta^\tp C^2\eta-\frac14\eta^\tp D_t\eta$.
\end{proof}

%======================================================================
\section{Main results}\label{sec:main}
%======================================================================

\begin{theorem}[Exact criterion at a fixed time]\label{thm:A}
Let $A\in\R^{d\times d}$, $B\in\R^{d\times m}$, let $C$ be symmetric positive definite, and let $t\ge0$.
\begin{enumerate}
\item[(a)] $K^tH_C\subset H_C$ if and only if $D_t\ge0$, i.e.
\begin{equation}\label{eq:critA}
e^{At}C^2e^{A^\tp t}\le C^2+4\Sigma(t).
\end{equation}
\item[(b)] If \eqref{eq:critA} holds, then $K^t$ is bounded on $H_C$ and $\|K^t\|_{H_C\to H_C}=e^{t\Tr(-A)/2}$.
\item[(c)] If \eqref{eq:critA} fails, there is $f\in H_C$ with $K^tf\notin H_C$, and $\hatf f$ can be chosen real,
even and supported in an open double cone.
\end{enumerate}
\end{theorem}

\begin{proof}
Let $f\in H_C$. By Lemma~\ref{lem:Kt}, $K^tf$ is real, continuous and in $L^2$. By
Lemma~\ref{lem:fourier}(b), $K^tf\in H_C$ if and only if the left-hand side of \eqref{eq:normid} is finite, and
then $\|K^tf\|_C^2=(2\pi)^{-d}\kappa_C^{-1}$ times that quantity.

\emph{(a), ``if'', and the upper bound in (b).} If $D_t\ge0$, then $\exp(-\frac14\eta^\tp D_t\eta)\le1$. By
\eqref{eq:normid} and \eqref{eq:normC}, $K^tf\in H_C$ and $\|K^tf\|_C^2\le e^{-t\Tr A}\|f\|_C^2$.

\emph{Lower bound in (b).} For $\varepsilon>0$ the indicator $F_\varepsilon:=\mathbf 1_{\{|\eta|<\varepsilon\}}$ is real,
even and belongs to $\FC{C}$, so $f_\varepsilon:=F_\varepsilon^\vee\in H_C\setminus\{0\}$ by Lemma~\ref{lem:fourier}(c). By
\eqref{eq:normid} and \eqref{eq:normC}, using $\eta^\tp D_t\eta\le\|D_t\|\varepsilon^2$ for $|\eta|<\varepsilon$,
\[
\|K^tf_\varepsilon\|_C^2=e^{-t\Tr A}(2\pi)^{-d}\kappa_C^{-1}\int_{|\eta|<\varepsilon}w_C(\eta)e^{-\frac14\eta^\tp D_t\eta}\dd\eta
\ \ge\ e^{-t\Tr A}e^{-\frac14\varepsilon^2\|D_t\|}\,\|f_\varepsilon\|_C^2 .
\]
Letting $\varepsilon\downarrow0$ gives $\|K^t\|^2\ge e^{-t\Tr A}$. Hence $\|K^t\|=e^{-t\Tr A/2}=e^{t\Tr(-A)/2}$.

\emph{(a), ``only if'', and (c).} Suppose $D_t\ge0$ fails. Lemma~\ref{lem:multiplier} with $S:=-D_t$ yields a real,
even $F_S\in\FC{C}$, supported in an open double cone, with $\int|F_S|^2w_Ce^{-\frac14\eta^\tp D_t\eta}\dd\eta=\infty$.
Let $f:=F_S^\vee$. Then $f\in H_C$ and $\hatf f=F_S$ by Lemma~\ref{lem:fourier}(c). By \eqref{eq:normid},
$\int|\hatf{K^tf}|^2w_C\dd\xi=\infty$, so $K^tf\notin H_C$.
\end{proof}

\begin{theorem}[Range of $K^t$]\label{thm:B}
For $t\ge0$ let
\begin{equation}\label{eq:Ctilde}
\Ct:=\Bigl(e^{-At}\bigl(C^2+4\Sigma(t)\bigr)e^{-A^\tp t}\Bigr)^{1/2},\qquad
\taut:=\frac{\det C}{\bigl(\det(C^2+4\Sigma(t))\bigr)^{1/2}} .
\end{equation}
Then $K^t$ maps $H_C$ bijectively onto $H_{\Ct}$, and $\|K^tf\|_{\Ct}^2=\taut\,\|f\|_C^2$ for all $f\in H_C$.
\end{theorem}

\begin{proof}
Let $f\in H_C$ and $g:=K^tf$, which is real, continuous and in $L^2$ by Lemma~\ref{lem:Kt}. By \eqref{eq:KtFT} and the
substitution $\xi=e^{A^\tp t}\eta$, exactly as in the proof of Proposition~\ref{prop:norm},
\[
\int|\hatf g(\xi)|^2w_{\Ct}(\xi)\dd\xi=e^{-t\Tr A}\int|\hatf f(\eta)|^2
\exp\Bigl(-\eta^\tp\Sigma(t)\eta+\tfrac14\eta^\tp e^{At}\Ct^{\,2}e^{A^\tp t}\eta\Bigr)\dd\eta
=e^{-t\Tr A}\int|\hatf f|^2w_C\dd\eta ,
\]
because $e^{At}\Ct^{\,2}e^{A^\tp t}=C^2+4\Sigma(t)$. Thus $g\in H_{\Ct}$ by Lemma~\ref{lem:fourier}(b) (with $\Ct$ in
place of $C$), and by \eqref{eq:normC}
\[
\|g\|_{\Ct}^2=\frac{\kappa_C}{\kappa_{\Ct}}\,e^{-t\Tr A}\,\|f\|_C^2=\frac{\det C}{\det\Ct}\,e^{-t\Tr A}\,\|f\|_C^2=\taut\,\|f\|_C^2,
\]
since $\det\Ct=e^{-t\Tr A}\det(C^2+4\Sigma(t))^{1/2}$. In particular $K^t|_{H_C}$ is injective.

For surjectivity let $g\in H_{\Ct}$ and define
$F(\eta):=e^{t\Tr A}\,\hatf g(e^{A^\tp t}\eta)\exp\bigl(\frac12\eta^\tp\Sigma(t)\eta\bigr)$. $F$ is Hermitian, and with
$\eta=e^{-A^\tp t}\xi$, $\mathrm d\eta=e^{-t\Tr A}\mathrm d\xi$, and
$\eta^\tp\Sigma(t)\eta+\frac14\eta^\tp C^2\eta=\frac14\xi^\tp\Ct^{\,2}\xi$,
\[
\int|F|^2w_C\dd\eta=e^{2t\Tr A}\int|\hatf g(e^{A^\tp t}\eta)|^2e^{\eta^\tp\Sigma(t)\eta+\frac14\eta^\tp C^2\eta}\dd\eta
=e^{t\Tr A}\int|\hatf g|^2w_{\Ct}\dd\xi<\infty .
\]
Hence $F\in\FC{C}$ and $f:=F^\vee\in H_C$ by Lemma~\ref{lem:fourier}(c). By \eqref{eq:KtFT},
$\hatf{K^tf}(\xi)=e^{-t\Tr A}F(e^{-A^\tp t}\xi)\exp(-\frac12\xi^\tp e^{-At}\Sigma(t)e^{-A^\tp t}\xi)=\hatf g(\xi)$. Since
$K^tf$ and $g$ are continuous $L^2$-functions with the same Fourier transform, $K^tf=g$.
\end{proof}

\begin{corollary}[On {\cite[Remark~4.2]{PSWPN}}]\label{cor:remark42}
For every $t\ge0$,
\[
K^tH_C\subset H_C\iff H_{\Ct}\subset H_C\iff \Ct^{\,2}\ge C^2\iff D_t\ge0 .
\]
Moreover $\Ct^{\,2}-C_t^2=2e^{-At}\Sigma(t)e^{-A^\tp t}\ge0$, so $H_{\Ct}\subset H_{C_t}$, and the inclusion is strict
whenever $\Sigma(t)\neq0$ (for instance for every $t>0$ if $(A,B)$ is controllable). The implication
``$K^tH_C\subset H_C\Rightarrow H_{C_t}\hookrightarrow H_C$'' considered in \cite[Remark~4.2]{PSWPN} is false in
general. For instance, in Example~\ref{ex:beta} with $4<\beta\le6$, $K^tH_C\subset H_C$ for all $t\ge0$, whereas
$C_t^2\ge C^2$ fails, and hence $H_{C_t}\not\subset H_C$, for all $t\in(0,\varepsilon)$ with some $\varepsilon>0$.
\end{corollary}

\begin{proof}
The first equivalence is Theorem~\ref{thm:B}, the second is Corollary~\ref{cor:inclusion}, and the third follows from
the congruence $\Ct^{\,2}-C^2=e^{-At}D_te^{-A^\tp t}$. The formula for $\Ct^{\,2}-C_t^2$ is immediate from
\eqref{eq:Ctilde} and \cite[(4.1)]{PSWPN}; Corollary~\ref{cor:inclusion} gives $H_{\Ct}\subset H_{C_t}$. If also
$H_{C_t}\subset H_{\Ct}$, then $C_t^2\ge\Ct^{\,2}$, i.e.\ $-2e^{-At}\Sigma(t)e^{-A^\tp t}\ge0$, which forces $\Sigma(t)=0$.
For the last assertion, let $f_x(t):=\ip{C_t^2x}x$ as in the proof of \cite[Theorem~2.1]{PSWPN}. There
$\dot f_x(0)=\ip{(2BB^\tp-(AC^2+C^2A^\tp))x}x$ is computed. If \eqref{eq:23} fails, some $x$ has $\dot f_x(0)<0$, so
$f_x(t)<f_x(0)=\|Cx\|^2$ for $t\in(0,\varepsilon)$, i.e.\ $C_t^2\ge C^2$ fails. By Corollary~\ref{cor:inclusion}, $H_{C_t}\not\subset H_C$
for these $t$. Invariance for all $t$ under \eqref{eq:23sharp} is Theorem~\ref{thm:C}.
\end{proof}

\begin{theorem}[Invariance for the whole semigroup]\label{thm:C}
Let $A\in\R^{d\times d}$, $B\in\R^{d\times m}$ and let $C$ be symmetric positive definite. The following are equivalent:
\begin{enumerate}
\item[(i)] $K^tH_C\subset H_C$ for every $t\ge0$;
\item[(ii)] there is a sequence $t_n\downarrow0$, $t_n>0$, with $K^{t_n}H_C\subset H_C$ for every $n$;
\item[(iii)] $\tfrac12(AC^2+C^2A^\tp)\le2BB^\tp$, i.e.\ \eqref{eq:23sharp}, i.e.\ $\Lambda_C\ge0$.
\end{enumerate}
If they hold, then $\|K^t\|_{H_C\to H_C}=e^{t\Tr(-A)/2}$ for every $t\ge0$.
\end{theorem}

\begin{proof}
We have $D_0=0$, and, since $A$ commutes with $e^{At}$,
\[
\frac{\mathrm d}{\mathrm dt}D_t=4e^{At}BB^\tp e^{A^\tp t}-Ae^{At}C^2e^{A^\tp t}-e^{At}C^2e^{A^\tp t}A^\tp=e^{At}\Lambda_Ce^{A^\tp t}.
\]
Hence
\begin{equation}\label{eq:Dint}
D_t=\int_0^te^{As}\Lambda_Ce^{A^\tp s}\dd s\qquad(t\ge0).
\end{equation}
(iii)$\Rightarrow$(i): if $\Lambda_C\ge0$, then each $e^{As}\Lambda_Ce^{A^\tp s}$ is positive semidefinite. So $D_t\ge0$
for all $t$, and Theorem~\ref{thm:A}(a) applies. (i)$\Rightarrow$(ii) is trivial. (ii)$\Rightarrow$(iii): by
Theorem~\ref{thm:A}(a), $t_n^{-1}D_{t_n}\ge0$. By \eqref{eq:Dint} and continuity of the integrand,
$t^{-1}D_t\to\Lambda_C$ as $t\downarrow0$. The cone of positive semidefinite matrices is closed, so $\Lambda_C\ge0$. The norm
statement is Theorem~\ref{thm:A}(b).
\end{proof}

\begin{remark}
Theorem~\ref{thm:C} is the computation in the proof of \cite[Theorem~2.1]{PSWPN} with $C_t$ replaced by $\Ct$. For
$\widetilde f_x(t):=\ip{\Ct^{\,2}x}x=\|Ce^{-A^\tp t}x\|^2+4\int_0^t\|B^\tp e^{-A^\tp s}x\|^2\dd s$ one finds
$\dot{\widetilde f}_x(t)=\ip{\Lambda_Ce^{-A^\tp t}x}{e^{-A^\tp t}x}$. Replacing $C_t$ by $\Ct$ is what turns the
sufficient condition into a characterisation.
\end{remark}

\begin{corollary}[Answer to the open problem]\label{cor:answer}
Condition \eqref{eq:23} implies \eqref{eq:23sharp}, but not conversely. Consequently \eqref{eq:23} is sufficient but
\emph{not necessary} for $K^tH_C\subset H_C$, $t\ge0$, even under all standing assumptions of \cite{PSWPN}. The
necessary and sufficient condition \eqref{eq:23sharp} is equivalent to condition \eqref{eq:23} for the matrix $C/\sqrt2$:
\[
K^tH_C\subset H_C\ \ \forall t\ge0\iff\tfrac12\Bigl(A\bigl(\tfrac{C}{\sqrt2}\bigr)^2+\bigl(\tfrac{C}{\sqrt2}\bigr)^2A^\tp\Bigr)\le BB^\tp .
\]
\end{corollary}

\begin{proof}
$BB^\tp\le2BB^\tp$ because $BB^\tp\ge0$. Examples~\ref{ex:beta} and~\ref{ex:scalar} satisfy \eqref{eq:23sharp} and violate
\eqref{eq:23}, and in both $A$ is Hurwitz and $(A,B)$ is controllable. The displayed reformulation holds because
$\frac12(A(C/\sqrt2)^2+(C/\sqrt2)^2A^\tp)=\frac14(AC^2+C^2A^\tp)$.
\end{proof}

\begin{corollary}[Isotropic kernels; sharp form of {\cite[Corollary~2.3]{PSWPN}}]\label{cor:iso}
For $\sigma>0$, $H_\sigma=H_{\sigma I}$ is invariant under every $K^t$, $t\ge0$, if and only if
$\frac12(A+A^\tp)\le\frac{2}{\sigma^2}BB^\tp$. In that case $\|K^t\|_{H_\sigma\to H_\sigma}=e^{t\Tr(-A)/2}$.
\end{corollary}

\begin{proof}
With $C=\sigma I$, \eqref{eq:23sharp} reads $\frac{\sigma^2}2(A+A^\tp)\le2BB^\tp$. Apply Theorem~\ref{thm:C}.
\end{proof}

%======================================================================
\section{Geometric form under the standing assumptions}\label{sec:geometry}
%======================================================================

\begin{theorem}[Invariant ellipsoids and the set of admissible times]\label{thm:D}
Let $A$ be Hurwitz (controllability is not needed), let $\Sigma=\int_0^\infty e^{As}BB^\tp e^{A^\tp s}\dd s$, and let
\[
P_C:=C^2+4\Sigma\quad(\text{positive definite}),\qquad |x|_{P_C}:=\bigl(x^\tp P_C^{-1}x\bigr)^{1/2},\qquad
\mathcal E_C:=\{x\in\R^d:\ |x|_{P_C}\le1\}.
\]
\begin{enumerate}
\item[(a)] $D_t=P_C-e^{At}P_Ce^{A^\tp t}$ for all $t\ge0$. Consequently, for each $t\ge0$,
\[
K^tH_C\subset H_C\iff e^{At}P_Ce^{A^\tp t}\le P_C\iff |e^{At}x|_{P_C}\le|x|_{P_C}\ \ \forall x\in\R^d\iff e^{At}\mathcal E_C\subset\mathcal E_C .
\]
\item[(b)] $A P_C+P_CA^\tp=-\Lambda_C$. Hence \eqref{eq:23sharp} holds iff $AP_C+P_CA^\tp\le0$ iff $\mathcal E_C$ is forward
invariant under $\dot x=Ax$. Likewise, \eqref{eq:23} holds iff
$A(C^2+2\Sigma)+(C^2+2\Sigma)A^\tp\le0$.
\item[(c)] The set $T_C:=\{t\ge0:\ K^tH_C\subset H_C\}$ is a closed additive subsemigroup of $[0,\infty)$ containing
$0$ and $[t_C,\infty)$ for every $t_C\ge0$ with $\sup_{t\ge t_C}\|e^{At}\|^2\le\lambda_{\min}(P_C)/\lambda_{\max}(P_C)$;
such $t_C$ exist. Either \eqref{eq:23sharp} holds and $T_C=[0,\infty)$, or $T_C\cap(0,\varepsilon)=\emptyset$ for some
$\varepsilon>0$.
\end{enumerate}
\end{theorem}

\begin{proof}
(a) Substituting $s=t+r$ in $\int_t^\infty e^{As}BB^\tp e^{A^\tp s}\dd s$ gives $\Sigma(t)=\Sigma-e^{At}\Sigma e^{A^\tp t}$.
Inserting this into \eqref{eq:Dt} yields $D_t=C^2+4\Sigma-e^{At}(C^2+4\Sigma)e^{A^\tp t}=P_C-e^{At}P_Ce^{A^\tp t}$, so the
first equivalence is Theorem~\ref{thm:A}(a). Write $P=P_C$, $E=e^{At}$ and $G:=P^{-1/2}EP^{1/2}$. Since
$P^{-1/2}EPE^\tp P^{-1/2}=GG^\tp$, we have $EPE^\tp\le P\iff GG^\tp\le I\iff\|G^\tp\|\le1\iff\|G\|\le1$. Moreover $\|G\|\le1$ means
$|P^{-1/2}EP^{1/2}z|\le|z|$ for all $z$, i.e.\ (with $x=P^{1/2}z$) $|Ex|_{P}\le|x|_{P}$ for all $x$. By homogeneity of
$|\cdot|_P$ this is equivalent to $E\mathcal E_C\subset\mathcal E_C$.

(b) As $A$ is Hurwitz, $A\Sigma+\Sigma A^\tp=\int_0^\infty\frac{\mathrm d}{\mathrm ds}\bigl(e^{As}BB^\tp e^{A^\tp s}\bigr)\dd s=-BB^\tp$.
Hence $AP_C+P_CA^\tp=AC^2+C^2A^\tp-4BB^\tp=-\Lambda_C$ and
$A(C^2+2\Sigma)+(C^2+2\Sigma)A^\tp=AC^2+C^2A^\tp-2BB^\tp$. By (a), forward invariance of $\mathcal E_C$ means
$e^{At}P_Ce^{A^\tp t}\le P_C$ for all $t\ge0$. Since
$\frac{\mathrm d}{\mathrm dt}e^{At}P_Ce^{A^\tp t}=e^{At}(AP_C+P_CA^\tp)e^{A^\tp t}$, this holds iff $AP_C+P_CA^\tp\le0$.
For ``only if'', divide $P_C-e^{At}P_Ce^{A^\tp t}\ge0$ by $t$ and let $t\downarrow0$.

(c) $0\in T_C$ because $D_0=0$. $T_C$ is closed because $t\mapsto D_t$ is continuous and the positive semidefinite cone is
closed. If $s,t\in T_C$, then, since congruence preserves the Loewner order,
$e^{A(s+t)}P_Ce^{A^\tp(s+t)}=e^{As}\bigl(e^{At}P_Ce^{A^\tp t}\bigr)e^{A^\tp s}\le e^{As}P_Ce^{A^\tp s}\le P_C$, so $s+t\in T_C$.
For $t\ge t_C$ we have $e^{At}P_Ce^{A^\tp t}\le\lambda_{\max}(P_C)\|e^{At}\|^2I\le\lambda_{\min}(P_C)I\le P_C$. Such $t_C$
exist because $\|e^{At}\|\to0$. The dichotomy is Theorem~\ref{thm:C}.
\end{proof}

\begin{remark}
Since $A\Sigma+\Sigma A^\tp=-BB^\tp\le0$, the ellipsoid of $\Sigma$ is always forward invariant. Theorem~\ref{thm:D}
states that $H_C$ is Koopman invariant precisely when adding $\frac14C^2$ to $\Sigma$ keeps the ellipsoid forward
invariant; \eqref{eq:23} demands this for $\frac12C^2$ instead of $\frac14C^2$. Theorem~\ref{thm:D}(c) also shows that for Hurwitz
$A$ \emph{every} Gaussian RKHS $H_C$ is invariant under $K^t$ for all sufficiently large $t$, even when invariance fails
for all small $t>0$ (Example~\ref{ex:beta}, $\beta>6$).
\end{remark}

%======================================================================
\section{Examples}\label{sec:examples}
%======================================================================

\begin{example}[A one-parameter family; $d=m=2$]\label{ex:beta}
Let $B=C=I_2$ and $A=A_\beta:=\begin{pmatrix}-1&\beta\\0&-1\end{pmatrix}$ with $\beta\ge0$. Then $A_\beta$ is Hurwitz (double
eigenvalue $-1$), and $(A_\beta,B)$ is controllable because $B$ is invertible. Here
$AC^2+C^2A^\tp=A_\beta+A_\beta^\tp=\begin{pmatrix}-2&\beta\\\beta&-2\end{pmatrix}$ has eigenvalues $-2\pm\beta$. Therefore
\[
\eqref{eq:23}\iff\tfrac12(\beta-2)\le1\iff\beta\le4,\qquad
\eqref{eq:23sharp}\iff\tfrac12(\beta-2)\le2\iff\beta\le6 .
\]
For $4<\beta\le6$, Theorem~\ref{thm:C} shows that $H_{I}$ is invariant under every $K^t$, $t\ge0$, with
$\|K^t\|_{H_I\to H_I}=e^{t}$ (as $\Tr(-A_\beta)=2$), although \eqref{eq:23} fails. For $\beta>6$ there is
$\varepsilon>0$ with $K^tH_I\not\subset H_I$ for all $t\in(0,\varepsilon)$.

Explicitly, $A_\beta=-I+\beta E_{12}$ with $E_{12}:=\begin{pmatrix}0&1\\0&0\end{pmatrix}$ and $E_{12}^2=0$, so
$e^{A_\beta t}=e^{-t}\begin{pmatrix}1&\beta t\\0&1\end{pmatrix}$ and
$\Sigma(t)=\int_0^te^{-2s}\begin{pmatrix}1+\beta^2s^2&\beta s\\\beta s&1\end{pmatrix}\dd s$. Using
$\int_0^te^{-2s}\dd s=\frac{1-e^{-2t}}2$, $\int_0^tse^{-2s}\dd s=\frac{1-e^{-2t}(1+2t)}4$ and
$\int_0^ts^2e^{-2s}\dd s=\frac{1-e^{-2t}(1+2t+2t^2)}4$, one obtains
\[
D_t=\begin{pmatrix}
3(1-e^{-2t})+\beta^2\bigl(1-e^{-2t}(1+2t+3t^2)\bigr) & \beta\bigl(1-e^{-2t}(1+3t)\bigr)\\[2pt]
\beta\bigl(1-e^{-2t}(1+3t)\bigr) & 3(1-e^{-2t})
\end{pmatrix},
\qquad D_t=t\begin{pmatrix}6&-\beta\\-\beta&6\end{pmatrix}+O(t^2)\ \ (t\downarrow0),
\]
in accordance with \eqref{eq:Dint}, since $\Lambda_C=4I-(A_\beta+A_\beta^\tp)=\begin{pmatrix}6&-\beta\\-\beta&6\end{pmatrix}$. Furthermore
$\Sigma=\begin{pmatrix}\frac12+\frac{\beta^2}4&\frac\beta4\\\frac\beta4&\frac12\end{pmatrix}$ and
$P_C=I+4\Sigma=\begin{pmatrix}3+\beta^2&\beta\\\beta&3\end{pmatrix}$. We find
$A_\beta P_C+P_CA_\beta^\tp=\begin{pmatrix}-6&\beta\\\beta&-6\end{pmatrix}$ and
$A_\beta(I+2\Sigma)+(I+2\Sigma)A_\beta^\tp=\begin{pmatrix}-4&\beta\\\beta&-4\end{pmatrix}$, which recovers the thresholds $6$ and
$4$ through Theorem~\ref{thm:D}(b).

\emph{Numerical illustration} (not used in any proof). For $\beta=8$, evaluating $\det D_t$ on a grid of
$(0,40]$ shows a single sign change, at $t_1\approx0.387611$ (bisection), with $D_t$ indefinite before and positive
semidefinite after it. For $t\ge40$ the bound $\|e^{A_8t}\|\le(1+8t)e^{-t}$ and the argument in the proof of
Theorem~\ref{thm:D}(c) give $D_t>0$. Thus, numerically, $T_{I}=\{0\}\cup[t_1,\infty)$ for $\beta=8$.
\end{example}

\begin{example}[Scalar noise; $d=2$, $m=1$]\label{ex:scalar}
Let $A=\begin{pmatrix}\frac32&-2\\2&-2\end{pmatrix}$, $B=\begin{pmatrix}1\\0\end{pmatrix}$ and $C=I_2$. Then $\Tr A=-\frac12$ and $\det A=1$, so
the eigenvalues $-\frac14\pm\frac{\sqrt{15}}4i$ have negative real part and $A$ is Hurwitz. Also
$\det[B,AB]=\det\begin{pmatrix}1&\frac32\\0&2\end{pmatrix}=2\neq0$, so $(A,B)$ is controllable. Since
$\frac12(A+A^\tp)=\operatorname{diag}(\frac32,-2)$ and $BB^\tp=\operatorname{diag}(1,0)$,
\[
BB^\tp-\tfrac12(A+A^\tp)=\operatorname{diag}\bigl(-\tfrac12,2\bigr)\not\ge0,\qquad
2BB^\tp-\tfrac12(A+A^\tp)=\operatorname{diag}\bigl(\tfrac12,2\bigr)\ge0 .
\]
Hence \eqref{eq:23} fails while \eqref{eq:23sharp} holds, and $H_I$ is invariant under every $K^t$ by Theorem~\ref{thm:C}.
\end{example}

%======================================================================
\section{Further remarks}\label{sec:remarks}
%======================================================================

\begin{lemma}[Inner products of Gaussian bumps]\label{lem:bumps}
Let $C',C''$ be symmetric positive definite and $y,y'\in\R^d$. Then $k^{C'}_y\in H_{C''}$ iff $G:=2C'^2-C''^2$ is
positive definite. In that case
\[
\ip{k^{C'}_y}{k^{C'}_{y'}}_{C''}=\frac{(\det C')^2}{\det C''\,(\det G)^{1/2}}\exp\bigl(-(y-y')^\tp G^{-1}(y-y')\bigr).
\]
\end{lemma}

\begin{proof}
By \eqref{eq:kernelFT}, $|\hatf{k^{C'}_y}|^2w_{C''}=\kappa_{C'}^2\exp(-\frac14\xi^\tp G\xi)$, which is integrable iff
$G>0$. If $G>0$, then by \eqref{eq:normC} and \eqref{eq:gauss} with $P=\frac14G$,
\[
\ip{k^{C'}_y}{k^{C'}_{y'}}_{C''}=(2\pi)^{-d}\frac{\kappa_{C'}^2}{\kappa_{C''}}\int e^{-\frac14\xi^\tp G\xi-i\xi\cdot(y-y')}\dd\xi
=(2\pi)^{-d}\frac{\pi^d(\det C')^2}{\pi^{d/2}\det C''}\,\pi^{d/2}2^d(\det G)^{-1/2}e^{-(y-y')^\tp G^{-1}(y-y')},
\]
and $(2\pi)^{-d}\pi^d2^d=1$.
\end{proof}

\begin{remark}[Why kernel expansions cannot detect non-invariance]\label{rem:sections}
By \cite[Proposition~4.1]{PSWPN}, $K^tk^C_z=\tau_t\,k^{C_t}_{e^{-At}z}$. By Lemma~\ref{lem:bumps}, $K^tk^C_z\in H_C$ iff
$2C_t^2-C^2>0$. Now $2C_t^2-C^2=e^{-At}(C^2+D_t)e^{-A^\tp t}$ and $D_t\to0$ as $t\downarrow0$. Hence for all
sufficiently small $t$, $K^t$ maps $\operatorname{span}\{k^C_z\}$ into $H_C$ \emph{regardless} of the sign of $D_t$. Non-invariance is
witnessed only by elements of $H_C$ whose Fourier transforms decay as slowly as membership in $H_C$ permits along
a cone of directions (Lemma~\ref{lem:multiplier}). Such functions are not finite kernel expansions. This explains why
the necessary condition $H_{0,C_t}\subset H_C$ noted in \cite[Remark~4.2]{PSWPN} is too weak to decide the question.
\end{remark}

\begin{remark}[Consistency with {\cite[Proposition~4.1]{PSWPN}}]\label{rem:consistency}
By Corollaries~\ref{cor:remark42} and~\ref{cor:inclusion}, $H_{\Ct}\subset H_{C_t}$ with
$\|g\|_{C_t}^2\le(\det\Ct/\det C_t)\|g\|_{\Ct}^2$. Together with Theorem~\ref{thm:B},
\[
\|K^tf\|_{C_t}^2\le\frac{\det\Ct}{\det C_t}\,\taut\,\|f\|_C^2
=\frac{\det(C^2+4\Sigma(t))^{1/2}}{\det(C^2+2\Sigma(t))^{1/2}}\cdot\frac{\det C}{\det(C^2+4\Sigma(t))^{1/2}}\,\|f\|_C^2=\tau_t\,\|f\|_C^2 .
\]
This is exactly the bound of \cite[Proposition~4.1]{PSWPN}. Theorem~\ref{thm:B} refines that proposition: the only
loss in the argument of \cite{PSWPN} is the inclusion $H_{\Ct}\subsetneq H_{C_t}$. On kernel sections Theorem~\ref{thm:B}
reads $\ip{K^tk^C_z}{K^tk^C_{z'}}_{\Ct}=\taut\,k^C(z,z')$. This also follows directly from
\cite[Proposition~4.1]{PSWPN} and Lemma~\ref{lem:bumps}, because $2C_t^2-\Ct^{\,2}=e^{-At}C^2e^{-A^\tp t}$.
\end{remark}

%======================================================================
\section{Numerical cross-checks}\label{sec:numerics}
%======================================================================

The script \texttt{verify.py} (NumPy/SciPy; fixed random seed) performs the checks below. They are floating-point
verifications of identities proved above, plus a Monte Carlo check of \cite[Proposition~4.1]{PSWPN}; no proof in this
note depends on them. $\Sigma(t)$ is computed by Van Loan's block exponential \cite{VanLoan} on steps of length
$\le0.25$, combined with the recursion $\Sigma(s+h)=\Sigma(h)+e^{Ah}\Sigma(s)e^{A^\tp h}$. The step size is needed
because the one-step formula cancels catastrophically for large $t$. All 34 checks passed.
\begin{enumerate}
\item $\Sigma(t)$ versus adaptive quadrature (random $A\in\R^{3\times3}$, $B\in\R^{3\times2}$, $t\in\{0.3,1.7,4\}$):
relative error $\le2.7\times10^{-13}$.
\item Identity \eqref{eq:Dint} for three random triples $(A,B,C)$, $A\in\R^{3\times3}$, $B\in\R^{3\times2}$, with $A$
\emph{not} Hurwitz (spectral abscissae $1.645$, $1.556$, $2.043$), $t\in\{0.05,0.8,2.5\}$:
relative error $\le5.2\times10^{-13}$; and $\|t^{-1}D_t-\Lambda_C\|/\max(1,\|\Lambda_C\|)\le3.7\times10^{-6}$ at $t=10^{-6}$,
consistent with the $O(t)$ remainder.
\item Theorem~\ref{thm:D}(a),(b) for a random Hurwitz $A\in\R^{3\times3}$, $B\in\R^{3\times1}$: $|D_t-(P_C-e^{At}P_Ce^{A^\tp t})|\le2.4\times10^{-14}$ for
$t\in\{0.1,1,6\}$ and $|AP_C+P_CA^\tp+\Lambda_C|\le1.8\times10^{-14}$.
\item Theorem~\ref{thm:B} on kernel sections (random $2\times2$ data, $t\in\{0.2,0.9\}$). The Gram matrix
$\bigl(\ip{K^tk^C_{z_i}}{K^tk^C_{z_j}}_{\Ct}\bigr)_{i,j\le4}$ is computed from \cite[Proposition~4.1]{PSWPN} and
Lemma~\ref{lem:bumps}; it agrees with $\taut\,(k^C(z_i,z_j))_{i,j}$ to within $3.8\times10^{-14}$. The Fourier-side value
$\|K^tk^C_z\|_C^2=e^{-t\Tr A}\det C\,\det(2C^2+4\Sigma(t)-e^{At}C^2e^{A^\tp t})^{-1/2}$, obtained from \eqref{eq:normid},
agrees with the value from \cite[Proposition~4.1]{PSWPN} to within $1.1\times10^{-13}$ (relative).
\item Monte Carlo check of \cite[Proposition~4.1]{PSWPN} (Example~\ref{ex:beta} with $\beta=5$, $t=0.4$, $2\times10^6$
exact Gaussian samples of $Z_t$): $0.436373\pm2.2\times10^{-4}$ versus the closed form $0.436051$ ($1.45$ standard
errors).
\item Example~\ref{ex:beta}: for $\beta\in\{4,5,6,7,8\}$, \eqref{eq:23} holds only for $\beta=4$ and \eqref{eq:23sharp}
holds exactly for $\beta\le6$. On a grid of $4200$ times in $(0,12]$, $D_t$ has no negative eigenvalue for $\beta\le6$.
For $\beta=7$ and $\beta=8$ the grid times at which $D_t$ is indefinite form an initial segment, ending at $0.298$
and $0.385$ respectively, and $D_t$ is positive semidefinite at all later grid times. For $\beta\in\{5,6\}$ the
matrix $C_t^2-C^2$ of \cite{PSWPN} has a negative eigenvalue ($-1.2\times10^{-2}$, resp.\ $-2.4\times10^{-2}$) at
some $t\le0.0124$, although $D_t\ge0$ there. The closed form of $D_t$ in Example~\ref{ex:beta} matches
the numerical $D_t$ to within $2.5\times10^{-14}$. For $\beta=8$, $\det D_t$ has exactly one sign change on $(0,40]$, at $t_1\approx0.387611$.
\item Example~\ref{ex:scalar}: eigenvalues $-0.25\pm0.968i$, $\operatorname{rank}[B,AB]=2$,
$\lambda_{\max}(\frac12(A+A^\tp)-BB^\tp)=+0.5$, $\lambda_{\max}(\frac12(A+A^\tp)-2BB^\tp)=-0.5$, and
$\min_t\lambda_{\min}(D_t)\ge0$ on a grid of $3000$ times in $(0,10]$.
\end{enumerate}

%======================================================================
\begin{thebibliography}{9}
%======================================================================

\bibitem{PSWPN}
F.~Philipp, M.~Schaller, K.~Worthmann, S.~Peitz, F.~N\"uske,
\emph{Invariance of Gaussian RKHSs under Koopman operators of stochastic differential equations with constant matrix
coefficients}, arXiv:2405.14429v1 (2024); published in Proc.\ Appl.\ Math.\ Mech.\ (PAMM),
doi:\href{https://doi.org/10.1002/pamm.202400127}{10.1002/pamm.202400127}.

\bibitem{Aronszajn}
N.~Aronszajn, \emph{Theory of reproducing kernels}, Trans.\ Amer.\ Math.\ Soc.\ \textbf{68} (1950), 337--404.

\bibitem{PaulsenRaghupathi}
V.~I.~Paulsen, M.~Raghupathi, \emph{An Introduction to the Theory of Reproducing Kernel Hilbert Spaces},
Cambridge Studies in Advanced Mathematics~152, Cambridge University Press, 2016.

\bibitem{SHS}
I.~Steinwart, D.~Hush, C.~Scovel, \emph{An explicit description of the reproducing kernel Hilbert spaces of Gaussian
RBF kernels}, IEEE Trans.\ Inform.\ Theory \textbf{52} (2006), 4635--4643.

\bibitem{VanLoan}
C.~F.~Van~Loan, \emph{Computing integrals involving the matrix exponential}, IEEE Trans.\ Automat.\ Control
\textbf{23} (1978), 395--404.

\bibitem{Wendland}
H.~Wendland, \emph{Scattered Data Approximation}, Cambridge Monographs on Applied and Computational Mathematics~17,
Cambridge University Press, 2005.

\end{thebibliography}

\end{document}
